\documentclass[10pt,a4paper]{amsart}
\usepackage{amsmath,amssymb,amsthm,graphicx}
% For commutative diagrams we use Xy-pic
\usepackage[all]{xy}

% We use 2cell for 2-commutative diagrams.
\xyoption{2cell}
% Theorem environments.
%
\theoremstyle{plain}
\newtheorem{theorem}[subsection]{Theorem}
\newtheorem{proposition}[subsection]{Proposition}
\newtheorem{lemma}[subsection]{Lemma}

\theoremstyle{definition}
\newtheorem{definition}[subsection]{Definition}
\newtheorem{example}[subsection]{Example}
\newtheorem{exercise}[subsection]{Exercise}
\newtheorem{situation}[subsection]{Situation}

\theoremstyle{remark}
\newtheorem{remark}[subsection]{Remark}
\newtheorem{remarks}[subsection]{Remarks}

\numberwithin{equation}{subsection}

% Macros
%
\def\lim{\mathop{\mathrm{lim}}\nolimits}
\def\colim{\mathop{\mathrm{colim}}\nolimits}
\def\Spec{\mathop{\mathrm{Spec}}}
\def\Hom{\mathop{\mathrm{Hom}}\nolimits}
\def\Ext{\mathop{\mathrm{Ext}}\nolimits}
\def\SheafHom{\mathop{\mathcal{H}\!\mathit{om}}\nolimits}
\def\SheafExt{\mathop{\mathcal{E}\!\mathit{xt}}\nolimits}
\def\Sch{\mathit{Sch}}
\def\Mor{\mathop{\mathrm{Mor}}\nolimits}
\def\Ob{\mathop{\mathrm{Ob}}\nolimits}
\def\Sh{\mathop{\mathit{Sh}}\nolimits}
\def\NL{\mathop{N\!L}\nolimits}
\def\CH{\mathop{\mathrm{CH}}\nolimits}
\def\proetale{{pro\text{-}\acute{e}tale}}
\def\etale{{\acute{e}tale}}
\def\QCoh{\mathit{QCoh}}
\def\Ker{\mathop{\mathrm{Ker}}}
\def\Im{\mathop{\mathrm{Im}}}
\def\Coker{\mathop{\mathrm{Coker}}}
\def\Coim{\mathop{\mathrm{Coim}}}

% Boxtimes
%
\DeclareMathSymbol{\boxtimes}{\mathbin}{AMSa}{"02}

%
% Macros for moduli stacks/spaces
%
\def\QCohstack{\mathcal{QC}\!\mathit{oh}}
\def\Cohstack{\mathcal{C}\!\mathit{oh}}
\def\Spacesstack{\mathcal{S}\!\mathit{paces}}
\def\Quotfunctor{\mathrm{Quot}}
\def\Hilbfunctor{\mathrm{Hilb}}
\def\Curvesstack{\mathcal{C}\!\mathit{urves}}
\def\Polarizedstack{\mathcal{P}\!\mathit{olarized}}
\def\Complexesstack{\mathcal{C}\!\mathit{omplexes}}
% \Pic is the operator that assigns to X its picard group, usage \Pic(X)
% \Picardstack_{X/B} denotes the Picard stack of X over B
% \Picardfunctor_{X/B} denotes the Picard functor of X over B
\def\Pic{\mathop{\mathrm{Pic}}\nolimits}
\def\Picardstack{\mathcal{P}\!\mathit{ic}}
\def\Picardfunctor{\mathrm{Pic}}
\def\Deformationcategory{\mathcal{D}\!\mathit{ef}}
\usepackage{cleveref}
\crefname{originalreference}{original reference}{original references}
\Crefname{originalreference}{Original reference}{Original references}
\makeatletter
\newlabel{item-surjective}{{06AP}{source}{}{sourcepointer-1}{}}
\newlabel{item-surjective@cref}{{[originalreference][0][]06AP}{[1][0][]source}{}{}{}}
\makeatother
\crefname{originalreference}{original reference}{original references}
\Crefname{originalreference}{Original reference}{Original references}
\makeatletter
\newlabel{item-vanishing}{{0A53}{source}{}{sourcepointer-2}{}}
\newlabel{item-vanishing@cref}{{[originalreference][0][]0A53}{[1][0][]source}{}{}{}}
\makeatother
\crefname{originalreference}{original reference}{original references}
\Crefname{originalreference}{Original reference}{Original references}
\makeatletter
\newlabel{lemma-ses-associated-to-open}{{095L}{source}{}{sourcepointer-3}{}}
\newlabel{lemma-ses-associated-to-open@cref}{{[originalreference][0][]095L}{[1][0][]source}{}{}{}}
\makeatother
\crefname{originalreference}{original reference}{original references}
\Crefname{originalreference}{Original reference}{Original references}
\makeatletter
\newlabel{proposition-finite-higher-direct-image-zero}{{03QP}{source}{}{sourcepointer-4}{}}
\newlabel{proposition-finite-higher-direct-image-zero@cref}{{[originalreference][0][]03QP}{[1][0][]source}{}{}{}}
\makeatother
\crefname{originallemma}{original lemma}{original lemmas}
\Crefname{originallemma}{Original lemma}{Original lemmas}
\makeatletter
\newlabel{sites-cohomology-lemma-apply-Leray}{{0733}{source}{}{sourcepointer-5}{}}
\newlabel{sites-cohomology-lemma-apply-Leray@cref}{{[originallemma][0][]0733}{[1][0][]source}{}{}{}}
\makeatother
\crefname{originalreference}{original reference}{original references}
\Crefname{originalreference}{Original reference}{Original references}
\makeatletter
\newlabel{situation-what-to-prove}{{0A52}{source}{}{sourcepointer-6}{}}
\newlabel{situation-what-to-prove@cref}{{[originalreference][0][]0A52}{[1][0][]source}{}{}{}}
\makeatother
\title{For a constructible sheaf the vanishing and surjectivity statements are unchanged by replacing it by its extension by zero from a dense open}
\author{The Stacks Project Authors}
\date{stacks-project:etale-cohomology.tex:local-snapshot-20261011; GNU FDL 1.2 or later}
\begin{document}
\maketitle
\noindent\textit{Editorial note.} All delivered mathematical text is the author's own, copied from the pinned Stacks Project chapter file etale-cohomology.tex; no statement or proof step was rewritten, reordered or removed.\par\medskip
\noindent\textit{Editorial note.} The counted proof is the author's complete argument: the proof is written out in full: the dimension count for the complement, the short exact sequence attached to the dense open, the identification of the quotient term through finiteness and Leray, and the conclusion from the vanishing on the zero-dimensional complement.\par\medskip
\noindent\textit{Editorial note.} The label lemma-restrict-to-open names the counted statement, and the proof environment immediately following it in the source is delivered verbatim as its proof.\par\medskip
\noindent\textit{Editorial note.} The counted statement is carried by the author's own lemma environment in the chapter source.\par\medskip
\noindent\textit{Editorial note.} The statement refers to the standing situation and to the enumerated statements by their labels; those references are kept as external pointers with their canonical Stacks tags.\par\medskip
\noindent\textit{Editorial note.} The preamble is the chapter's own native preamble excerpt already used for this source: the Xy-pic setup (native preamble lines 14-18) and the native theorem declarations with the mathematical macros (native preamble lines 154-217).\par\medskip
\noindent\textit{Editorial note.} The same-name label in the chapter chow.tex was compared with the counted statement and found to be a different proposition (Chow groups of a closed subscheme), so no cross-chapter restriction applies.\par\medskip
\noindent\textit{Allowed standard prerequisites (editorial).}\begin{itemize}
\item The short exact sequence associated to a dense open subscheme. Same chapter, the short exact sequence 0 -> j\_!j\textasciicircum{}\{-1\}F -> F -> i\_*i\textasciicircum{}\{-1\}F -> 0 for the inclusion j of a dense open with closed complement i. Supplies the exact sequence whose long exact cohomology sequence drives the comparison.
\item Vanishing of the higher direct images of a finite morphism and the Leray spectral sequence. Same chapter, the vanishing of R\textasciicircum{}q i\_* for the finite closed immersion i : Z -> X and the Leray spectral sequence, whose combination gives H\textasciicircum{}q(X, i\_*i\textasciicircum{}\{-1\}F) = H\textasciicircum{}q(Z, i\textasciicircum{}*F). Supplies the identification of the cohomology of the third term with the cohomology of the zero-dimensional set Z.
\item The dimension zero structure of the complement. Same chapter, the standing situation in which the complement of a dense open in the curve is a finite set of k-rational points. Supplies that Z is a disjoint union of spectra of algebraically closed fields, hence has no higher cohomology.
\end{itemize}\medskip
\ifdefined\hypertarget\hypertarget{sourcepointer-1}{}\fi
\noindent\textit{Source reference.} Original reference 06AP: sources/stacks-etale-cohomology/pinned-original-source.tex, source lines 15365--15365. Prerequisite pointer invoked at proof line(s) 15495; it names a standard fact of the same chapter that the counted argument uses, and every step of the counted proof built on it is retained verbatim in the delivered body.\par\medskip
\ifdefined\hypertarget\hypertarget{sourcepointer-2}{}\fi
\noindent\textit{Source reference.} Original reference 0A53: sources/stacks-etale-cohomology/pinned-original-source.tex, source lines 15335--15335. Prerequisite pointer invoked at proof line(s) 15495; it names a standard fact of the same chapter that the counted argument uses, and every step of the counted proof built on it is retained verbatim in the delivered body.\par\medskip
\ifdefined\hypertarget\hypertarget{sourcepointer-3}{}\fi
\noindent\textit{Source reference.} Original reference 095L: sources/stacks-etale-cohomology/pinned-original-source.tex, source lines 11588--11588. Prerequisite pointer invoked at proof line(s) 15506; it names a standard fact of the same chapter that the counted argument uses, and every step of the counted proof built on it is retained verbatim in the delivered body.\par\medskip
\ifdefined\hypertarget\hypertarget{sourcepointer-4}{}\fi
\noindent\textit{Source reference.} Original reference 03QP: sources/stacks-etale-cohomology/pinned-original-source.tex, source lines 8245--8245. Prerequisite pointer invoked at proof line(s) 15510; it names a standard fact of the same chapter that the counted argument uses, and every step of the counted proof built on it is retained verbatim in the delivered body.\par\medskip
\ifdefined\hypertarget\hypertarget{sourcepointer-5}{}\fi
\noindent\textit{Source reference.} Original lemma 0733: sources/stacks-chow/tags.txt, source lines 9047--9047. Prerequisite pointer invoked at proof line(s) 15512; it names a standard fact of a sibling chapter that the counted argument uses, and every step of the counted proof built on it is retained verbatim in the delivered body.\par\medskip
\ifdefined\hypertarget\hypertarget{sourcepointer-6}{}\fi
\noindent\textit{Source reference.} Original reference 0A52: sources/stacks-etale-cohomology/pinned-original-source.tex, source lines 15324--15324. Prerequisite pointer invoked at proof line(s) 15492; it names a standard fact of the same chapter that the counted argument uses, and every step of the counted proof built on it is retained verbatim in the delivered body.\par\medskip
\begin{lemma}
\label{lemma-restrict-to-open}
In Situation \ref{situation-what-to-prove} assume $\mathcal{F}$
constructible. Let $j : X' \to X$ be the inclusion of a dense open subscheme.
Then statements
(\ref{item-vanishing}) -- (\ref{item-surjective}) hold for $\mathcal{F}$
if and only if they hold for $j_!j^{-1}\mathcal{F}$.
\end{lemma}

\begin{proof}
Since $X'$ is dense, we see that $Z = X \setminus X'$ has dimension $0$
and hence is a finite set $Z = \{x_1, \ldots, x_n\}$ of $k$-rational points.
Consider the short exact sequence
$$
0 \to j_!j^{-1}\mathcal{F} \to \mathcal{F} \to i_*i^{-1}\mathcal{F} \to 0
$$
of Lemma \ref{lemma-ses-associated-to-open}. Observe that
$H^q_\etale(X, i_*i^{-1}\mathcal{F}) = H^q_\etale(Z, i^*\mathcal{F})$.
Namely, $i : Z \to X$ is a closed immersion, hence finite,
hence we have the vanishing of $R^qi_*$ for $q > 0$ by
Proposition \ref{proposition-finite-higher-direct-image-zero}, and
hence the equality follows from the Leray spectral sequence
(Cohomology on Sites, Lemma \ref{sites-cohomology-lemma-apply-Leray}).
Since $Z$ is a disjoint union of spectra of algebraically closed
fields, we conclude that $H^q_\etale(Z, i^*\mathcal{F}) = 0$
for $q > 0$ and
$$
H^0_\etale(Z, i^{-1}\mathcal{F}) =
\bigoplus\nolimits_{i = 1, \ldots, n} \mathcal{F}_{x_i}
$$
which is finite as $\mathcal{F}_{x_i}$ is finite due to the
assumption that $\mathcal{F}$ is constructible.
The long exact cohomology sequence gives an exact sequence
$$
0 \to
H^0_\etale(X, j_!j^{-1}\mathcal{F}) \to
H^0_\etale(X, \mathcal{F}) \to
H^0_\etale(Z, i^{-1}\mathcal{F}) \to
H^1_\etale(X, j_!j^{-1}\mathcal{F}) \to
H^1_\etale(X, \mathcal{F}) \to 0
$$
and isomorphisms $H^q_\etale(X, j_!j^{-1}\mathcal{F}) \to
H^q_\etale(X, \mathcal{F})$ for $q > 1$.

\medskip\noindent
At this point it is easy to deduce each of 
(\ref{item-vanishing}) -- (\ref{item-surjective})
holds for $\mathcal{F}$ if and only if it holds for
$j_!j^{-1}\mathcal{F}$. We make a few small remarks to help the reader:
(a) if $\mathcal{F}$ is torsion prime to the characteristic
of $k$, then so is $j_!j^{-1}\mathcal{F}$,
(b) the sheaf $j_!j^{-1}\mathcal{F}$ is constructible,
(c) we have $H^0_\etale(Z, i^{-1}\mathcal{F}) =
H^0_\etale(Z_{k'}, i^{-1}\mathcal{F}|_{Z_{k'}})$, and (d)
if $U \subset X$ is an open, then $U' = U \cap X'$ is dense in $U$.
\end{proof}

\end{document}
